\textbf{AGENTE INTELIGENTE}: entidad de software capaz de percibir su entorno, razonar sobre la información disponible y ejecutar acciones autónomas para alcanzar un objetivo definido, utilizando herramientas externas cuando es necesario.\\

\textbf{CHAIN-OF-THOUGHT (CoT)}: técnica de prompting que solicita al modelo de lenguaje descomponer su razonamiento en pasos intermedios explícitos antes de emitir una respuesta final.\\

\textbf{DESINFORMACIÓN}: información falsa o engañosa creada y difundida deliberadamente con el propósito de causar daño o confundir a la audiencia.\\

\textbf{FACT-CHECKING}: proceso de verificación de hechos que consiste en evaluar la veracidad de afirmaciones, noticias o declaraciones mediante la consulta de fuentes primarias y evidencia verificable.\\

\textbf{FEW-SHOT PROMPTING}: técnica en la que se proporcionan al modelo de lenguaje un número reducido de ejemplos (típicamente entre 3 y 5) de la tarea deseada antes de presentarle la instancia a resolver.\\

\textbf{IFCN}: International Fact-Checking Network, red internacional que certifica organizaciones de verificación de hechos bajo principios de transparencia, imparcialidad y metodología rigurosa.\\

\textbf{LANGGRAPH}: biblioteca de orquestación de agentes basada en grafos de estados que permite definir flujos de decisión complejos con nodos condicionales y herramientas externas.\\

\textbf{LLM}: Large Language Model o modelo de lenguaje de gran escala. Modelo de inteligencia artificial entrenado sobre grandes volúmenes de texto capaz de comprender y generar lenguaje natural.\\

\textbf{PROMPTING}: técnica de interacción con modelos de lenguaje mediante instrucciones en lenguaje natural que definen la tarea, el contexto, el formato de salida y las restricciones del proceso.\\

\textbf{REACT}: framework propuesto por Yao et al. (2022) que combina razonamiento (Reasoning) y acción (Acting) en modelos de lenguaje, permitiendo que el modelo genere trazas de pensamiento intercaladas con acciones sobre herramientas externas.\\

\textbf{ZERO-SHOT}: configuración en la que el modelo de lenguaje recibe únicamente la instrucción de la tarea sin ejemplos previos, dependiendo exclusivamente de su conocimiento preentrenado.
